\documentclass[12pt]{article}
\usepackage[utf8]{inputenc}
\usepackage[T1]{fontenc}
\usepackage{hyperref}
\title{Kula Language Resources for K-12 Curriculum Development}
\author{Compiled by Claude Sonnet 5 (claude-sonnet-5) for Andr\'as Kornai}
\date{}
\begin{document}
\maketitle
\section{Introduction}
Kula (ISO 639-3 tpg; Glottocode kula1280) is a Timor-Alor-Pantar language of Alor Island, Indonesia. It is also called Lamtoka or Tanglapui. It has about 5,000 speakers. Searches found only a little: one reference grammar, some phonology work and an archive of documentation material. They found no Kula school materials, no primers and no dictionary.
\section{Primary Educational Resources}
No primers, textbooks or school materials were found in search results. Kula is not known to be taught in schools, and one source says its use is discouraged there.
\section{Linguistic and Reference Materials}
\begin{itemize}
\item Wikipedia article on the Kula language, which cites Williams (2017), ``Kula'', in Schapper (ed.), \emph{The Papuan Languages of Timor, Alor and Pantar}, vol.\ 2 (De Gruyter): \url{https://en.wikipedia.org/wiki/Kula_language}
\item Glottolog entry: \url{https://glottolog.org/resource/languoid/id/kula1280}
\item New Guinea World page for the Alor-Tip group, with a bibliography that includes Haan (2006) on Kula phonology: \url{https://newguineaworld.linguistik.uzh.ch/families/trans-new-guinea/berau-gulf/west-bomberai/timor-alor-pantar/alor-pantar/east-alor/alor-tip/start}
\end{itemize}
\section{Religious and Community Resources}
None were found in search results.
\section{Digital and Audio Resources}
\begin{itemize}
\item Endangered Languages Project sample page for Kula: \url{https://www.endangeredlanguages.com/lang/2086/samples/11490}
\end{itemize}
The Wikipedia article links an ELAR archive, ``Documenting Language and Interaction in Kula''. I did not see its own URL in the results.
\section{Curriculum Development Context}
Kula is endangered. Speakers have been resettled, and school use is discouraged. A curriculum would have to start from the reference grammar and the documentation archive, working with the community. Dictionary and primer materials would have to be created. The published Sawila dictionary (Kratochvil et al., 2014) covers a related language only.
\end{document}